\section{总结与展望}
\label{sec:conclusion_v2}

\subsection{论文主要工作}

本文以燃煤机组汽轮机热力系统为研究对象，针对宽负荷运行条件下正常状态基准难以确定、设备与故障知识分散以及高质量故障样本有限等问题，围绕设备健康状态建模、汽轮机故障知识图谱构建和KG-GCN故障诊断方法开展研究，形成了由动态运行数据与静态领域知识共同驱动的状态监测与故障诊断流程。主要研究工作和结论如下。

（1）针对汽轮机状态参数随负荷、蒸汽边界和运行调节持续变化的问题，建立了基于设备实际物理边界的GRU健康状态监测方法。按照$U+B\rightarrow Y$的设备级建模方式，将操作变量及上游/外部热力边界作为输入，以设备本体和热力接口状态作为输出，使模型学习给定运行条件下的正常响应关系。选取VHP作为代表性对象，利用2024年5月正常运行DCS数据训练，并采用2024年6月数据进行独立测试。ANN、TCN、iTransformer-small和GRU对比结果表明，GRU在最终24个状态参数上的宏平均$R^2$达到0.9511，在VHP本体、一级抽汽及一次冷再接口等主要状态上均获得较好的预测效果。由此建立的动态健康参考能够削弱正常工况变化对状态评价的干扰，并为后续利用实测值与健康预测值之间的偏差构造故障诊断动态特征提供基础。

（2）针对汽轮机设备连接关系复杂以及故障知识分散的问题，构建了基于运行生产资料的汽轮机系统故障知识图谱。结合本机DCS点表、设备输入输出关系和系统接口资料确定设备、参数及热力拓扑，并利用同行评议论文、标准和公开运行事件补充具有通用物理意义的故障机理。根据故障诊断需求定义设备、参数、异常、故障和处置等实体及相应关系类型，利用大语言模型完成候选三元组抽取，再结合实体类型、关系方向、设备兼容性、来源证据和Embedding语义对齐进行知识校核和实体标准化。当前审计知识库共包含127个标准实体和123条候选关系，其中121条关系进入主知识图谱，2条存在资料差异的候选关系保留待核。该知识图谱实现了汽轮机主通流、再热、抽汽回热和凝汽冷端等设备结构知识与故障机理知识的统一图结构表达。

（3）针对不同故障共享局部监测参数时容易产生类别混淆以及故障标签数量有限的问题，提出了基于知识图谱与图卷积网络的KG-GCN故障诊断方法。该方法首先将健康状态模型得到的多参数残差组织为滑动时间窗口，并映射为知识图谱参数节点的动态特征；随后将知识图谱投影得到的参数关系作为图结构输入，通过GCN沿设备和参数关联关系进行逐层特征聚合。训练阶段先在不使用类别标签的训练残差窗口上进行随机遮蔽重构预训练，再固定图卷积编码器并利用带标签故障窗口训练分类层。以VHP为代表性对象，基于真实健康运行残差与独立故障机理构造正常状态、VHP通流部件侵蚀、一级抽汽支路泄漏和VHP汽封泄漏四类半物理诊断场景。固定随机种子主实验中，KG-GCN的Accuracy和F1分别达到98.57\%和98.58\%，高于CNN和AE；3组随机种子重复实验的平均Macro-F1为98.46\%$\pm$0.11\%。在分类阶段仅使用20\%带标签故障窗口时，掩码预训练使Macro-F1由94.48\%提高至96.67\%，说明该训练方式能够在故障标签有限条件下改善模型诊断性能与结果稳定性。

综上，本文形成了“设备健康状态建模—状态残差构建—汽轮机故障知识图谱—图卷积结构化特征聚合—掩码预训练—故障分类”的完整技术流程。其中，知识图谱及KG-GCN方法面向燃煤机组汽轮机热力系统构建，VHP用于代表性算例验证。当前结果表明，健康状态残差能够为变工况下故障诊断提供统一动态表征，知识图谱能够为多参数残差提供实际热力结构约束，二者融合为汽轮机复杂故障诊断提供了一条数据与知识协同的技术路径。

\subsection{后续研究方向}

本文工作仍存在以下需要进一步研究的问题。

（1）\textbf{现场真实故障样本验证。} 当前故障分类实验采用真实健康运行残差背景，并依据独立故障机理构造半物理故障场景，能够用于验证模型在受控条件下的分类能力，但不能替代现场真实故障事件验证。后续应进一步收集机组缺陷记录、检修记录和历史故障数据，对不同故障的发生时刻、严重程度及多参数动态响应进行校准；在现场故障数据仍不足时，可结合经实际DCS数据校准的热力仿真模型补充不同故障程度和动态演化过程下的样本。

（2）\textbf{多设备及跨设备故障扩展。} 本文知识图谱覆盖汽轮机主通流、再热、抽汽回热及凝汽冷端等主要对象，但当前故障诊断算例仅对VHP任务子图进行了完整实验。后续可按照相同的设备级健康建模、参数节点映射和KG-GCN训练流程扩展至HP、IP、LP、凝汽器及回热设备，并进一步研究故障在相邻设备之间的传播、复合故障耦合以及系统级多故障并发识别问题。

（3）\textbf{知识图谱增量更新与不确定性建模。} 随着机组运行、检修和技术资料持续积累，汽轮机故障知识图谱需要具备增量更新能力。后续可建立面向运行日志、检修报告、技术通报和专家经验的候选知识抽取与审核机制，并进一步研究关系证据强度、边权重、动态图结构以及知识不确定性对GCN特征传播和故障诊断结果的影响，使领域知识能够随设备状态和运行经验持续更新。

（4）\textbf{复杂运行条件下的工程鲁棒性。} 实际在线应用还可能面临测点缺失、传感器漂移、工况分布变化和长期设备性能退化等问题。后续可将测点可靠性、缺失数据重构和在线分布漂移检测与KG-GCN诊断框架结合，并通过长期连续运行数据评价模型在不同季节、不同检修周期和不同负荷分布下的稳定性，为最终工程部署提供更加充分的验证依据。
